\chapter{Treasury and Funding}\label{ch:fm:treasury-and-funding}

In March 2020 initial margin requirements at the world's clearing houses rose by roughly \$300 billion in a month, and the variation margin called
in both cleared and bilateral markets was far larger than in February. Every firm with cleared or margined positions had to find the cash the
morning each call came, from wherever its treasury had left it: at a broker that would release it, in a money market fund that could be sold that
day, or in a credit line that the bank would honour. A firm that had kept its cash where it could not reach it learned its \omterm{def:fm:treasury-and-funding:horizon}{survival horizon} the
hard way.

\section{What treasury does in a trading firm}

A trading firm's treasury holds and moves its cash and collateral. It keeps enough at each prime broker, clearing broker and clearing house to meet
their requirements; it finances long positions and borrows securities for short ones at the lowest cost; it forecasts what the next days will
demand; and it keeps a reserve for the days when the forecast is wrong. It is the function whose failure is fastest: a firm with good positions and
no cash for a margin call is in default by the evening.

\begin{definition}[Unencumbered cash, liquidity buffer]\label{def:fm:treasury-and-funding:cash}
\emph{Unencumbered cash}\index{unencumbered cash} is cash, and assets convertible to cash the same day, that the firm holds free of any pledge,
margin requirement or restriction and can use at once. A \emph{liquidity buffer}\index{liquidity buffer} is the part of it held against stressed
outflows (margin calls, collateral returns, redemptions, loss of financing) over a defined horizon, and not used to fund positions.
\end{definition}

\begin{omfigure}
\begin{tikzpicture}[font=\footnotesize,>=stealth,box/.style={draw,rounded corners,align=center,minimum height=0.75cm,text width=2.4cm}]
\node[box,fill=omDef!15,text width=3cm] (t) at (0,0) {treasury\\{\scriptsize unencumbered cash}};
\node[box,fill=omProp!15] (a) at (-4.2,1.9) {prime broker A};
\node[box,fill=omProp!15] (b) at (0,2.3) {prime broker B};
\node[box,fill=omProp!15] (c) at (4.2,1.9) {prime broker C};
\node[box,fill=omThm!15] (f) at (4.6,-2.3) {clearing broker\\{\scriptsize futures}};
\node[box,fill=gray!15] (ccp) at (4.6,-3.8) {clearing house};
\node[box,fill=gray!15] (mmf) at (-4.6,-2.3) {money market fund, bills};
\node[box,fill=gray!15] (rcf) at (0,-2.8) {committed credit line};
\draw[<->] (t) -- (a); \node[font=\scriptsize] at (-2.9,0.75) {calls, releases}; \draw[<->] (t) -- (b); \draw[<->] (t) -- (c);
\draw[<->] (t) -- node[pos=0.5,above,sloped,font=\scriptsize] {variation margin} (f); \draw[<->] (f) -- (ccp);
\draw[<->] (t) -- node[pos=0.5,above,sloped,font=\scriptsize] {same-day sale} (mmf); \draw[->] (rcf) -- node[right,font=\scriptsize] {draw} (t);
\end{tikzpicture}

\omcaption{A trading firm's treasury: the accounts that hold its collateral and call or release margin daily, and the sources it can turn into cash
the same day. Each arrow has its own terms, cut-off times and notice periods. Schematic.}\label{fig:fm:treasury-and-funding:map}
\end{omfigure}

The bank version of the idea is Book 6's liquidity coverage ratio, and the internal price of using the balance sheet is its funds-transfer pricing.
A proprietary firm or a fund has no regulatory ratio to meet in most places; it has its brokers' terms, its own positions and the market's moves.

\section{Prime-broker relationships and margin terms}

A prime broker (Book 1) finances the firm's longs, lends it the securities it shorts, holds its positions and collateral, and sets the margin it
requires on them. The regulatory minimum for US customer accounts is only the floor (\cref{dat:fm:treasury-and-funding:rules}); what binds a
professional client is the broker's own rule.

\begin{definition}[House margin]\label{def:fm:treasury-and-funding:house}
\emph{House margin}\index{house margin} is the margin a broker requires of a client under its own rules, above any regulatory or clearing-house
minimum: typically a rate on gross exposure, a rate on the net, add-ons for concentration, liquidity and correlation, and a methodology the broker
may change at its discretion unless the agreement says otherwise.
\end{definition}

\begin{definition}[Margin lock-up]\label{def:fm:treasury-and-funding:lockup}
A \emph{margin lock-up}\index{margin lock-up} is a broker's contractual commitment not to change its margin terms (rates, methodology, eligible
collateral) or withdraw its financing for a client without a notice period, commonly thirty to ninety days, except on defined events such as a
large fall in the client's net asset value.
\end{definition}

\begin{dated}{2026-09}{Regulatory margin floors for US customer accounts}\label{dat:fm:treasury-and-funding:rules}
Regulation T (12 CFR 220.12) sets the initial margin on a margin equity security at 50\% of its current market value, and on a short sale of a
non-exempt security at 150\% of its value. FINRA Rule 4210 sets a maintenance margin of 25\% of the current market value of long margin securities,
and lets member firms apply risk-based portfolio margin instead of the strategy-based rules. Brokers' \omterm{def:fm:treasury-and-funding:house}{house margins} sit above these floors; offshore
and institutional arrangements are set by contract.
\end{dated}

Three brokers with different rules price the same book very differently. The chapter's book is long \$390 million and short \$330 million in ten
equity positions, hedged with a \$60 million short index future at a clearing broker (illustrative). Broker A charges 4\% of gross plus 25\% of
the absolute net, with a 15\% add-on on any position above \$60 million; broker B charges a flat 10\% of gross; broker C charges 5\% of gross plus
15\% of net, with a 30\% add-on above \$40 million. A single \$120 million long costs \$43.8 million of margin at A, \$12.0 million at B and \$48.0
million at C: A and C reward a balanced book and punish concentration, B charges for size alone.

\omcode{code/firm/treasury/firm_treasury.py}{35}{55}{House-margin rules as data: gross, net and concentration terms, and the financing spreads on
longs and shorts.}\label{lst:fm:treasury-and-funding:rules}

\section{Margin optimisation}

\begin{definition}[Margin optimisation]\label{def:fm:treasury-and-funding:optim}
\emph{Margin optimisation}\index{margin optimisation} is the choice of where to hold each position, and how to pair positions within each account,
so that the total margin and financing cost across a firm's brokers and clearing houses is least for the same book.
\end{definition}

\begin{proposition}[Optimising never costs more than one broker]\label{prop:fm:treasury-and-funding:lp}
If each broker's margin is a sum of rates times gross exposure, absolute net exposure and excesses of positions over caps, and financing is linear
in the positions, the annual cost is a convex, piecewise-linear function of the fractions $x_{ib}\ge0$ of each position $i$ held at broker $b$,
$\sum_bx_{ib}=1$. Its minimum is the solution of a linear programme, and it is no larger than the cost of holding the whole book at any single
broker.
\end{proposition}

\begin{proof}
Gross exposure is linear in $x$; the absolute value of a linear function and the positive part of an affine one are convex; sums of convex
functions with non-negative weights are convex. Replacing each absolute value by a variable $t_b\ge\pm\sum_ip_ix_{ib}$ and each positive part by
$u_{ib}\ge|p_i|x_{ib}-\text{cap}_b$, $u_{ib}\ge0$, gives a linear programme with the same minimum. Holding everything at one broker is a
feasible point, so the minimum is no larger than its cost.
\end{proof}

\begin{table}[ht]
\centering\small
\begin{tabular}{lrrrrr}
\toprule
allocation & margin A & margin B & margin C & total margin & annual cost \\
\midrule
all at A & 71.55 & -- & -- & 71.55 & 6.88 \\
all at B & -- & 72.00 & -- & 72.00 & 5.99 \\
all at C & -- & -- & 144.00 & 144.00 & 10.92 \\
pro rata & 14.60 & 24.00 & 15.00 & 53.60 & 5.82 \\
optimised & 20.40 & 6.00 & 7.50 & 33.90 & 4.98 \\
\bottomrule
\end{tabular}
\caption{\omterm{def:fm:treasury-and-funding:house}{House margin} and annual cost (margin at a 5\% cost of cash plus financing) of the chapter's book under five allocations across three
prime brokers, \$ million. Data: \texttt{fm\_treasury.summary}.}\label{tab:fm:treasury-and-funding:alloc}
\end{table}

The optimiser holds each broker's slice balanced, so that broker A's net charge and C's concentration add-ons almost vanish, and sends B only what
its flat gross rate prices cheaply (\cref{tab:fm:treasury-and-funding:alloc}, \cref{fig:fm:treasury-and-funding:margins}). Margin falls from \$72.0
million at the cheapest single broker to \$33.9 million, 53\% less, and the annual cost from \$5.99 million to \$4.98 million. Splitting the book
equally goes about half the way, to \$53.6 million.

\begin{omfigure}
\begin{tikzpicture}
\begin{axis}[width=10.5cm,height=5cm,ybar stacked,bar width=18pt,ylabel={\small house margin (\$ million)},ymin=0,ymax=150,xtick={0,1,2,3,4},
  xticklabels={all at A,all at B,all at C,pro rata,optimised},xmin=-0.6,xmax=4.6,
  tick label style={font=\footnotesize},grid=major,grid style={gray!20},table/col sep=comma,
  legend cell align=left,legend style={font=\scriptsize,at={(0.5,-0.25)},anchor=north,/tikz/every even column/.append style={column sep=8pt}},legend columns=3]
\addplot[fill=omDef!55,draw=omDef] table[x=pos,y=A]{figdata/desk/14-treasury-and-funding/margins.csv};
\addplot[fill=omProp!55,draw=omProp] table[x=pos,y=B]{figdata/desk/14-treasury-and-funding/margins.csv};
\addplot[fill=omThm!50,draw=omThm] table[x=pos,y=C]{figdata/desk/14-treasury-and-funding/margins.csv};
\legend{broker A,broker B,broker C}
\end{axis}
\end{tikzpicture}

\omcaption{\omterm{def:fm:treasury-and-funding:house}{House margin} on the same book under five allocations across three prime brokers. The linear programme of
\cref{prop:fm:treasury-and-funding:lp} keeps each broker's slice balanced and cuts margin by 53\% against the cheapest single broker. Data:
\texttt{fm\_treasury.summary}.}\label{fig:fm:treasury-and-funding:margins}
\end{omfigure}

\begin{remark}[What the optimiser does not see]\label{rem:fm:treasury-and-funding:limits}
Brokers price relationships, not slices: a broker that sees only the balanced remainder of a book may reprice it, and one that holds little of it
has little reason to support the firm in stress. Moving positions between brokers costs money and time, and in stress may not be possible. The
optimiser gives the cheapest allocation for today's terms; the treasurer chooses one that remains acceptable when the terms change.
\end{remark}

\section{Forecasting cash and collateral}

A cash forecast follows each account through the days ahead. Each account holds collateral; each day it gains or loses the P\&L of the positions it
holds and faces a requirement; a shortfall is a call paid from \omterm{def:fm:treasury-and-funding:cash}{unencumbered cash} the same day, and an excess is returned to cash if the broker
releases it (\cref{lst:fm:treasury-and-funding:forecast}). The clearing broker's requirement on the future comes from firm.initmargin's filtered
historical simulation (Book 6), and the prime brokers' from their house rules times a multiplier they may raise in stress.

\omcode{code/firm/treasury/firm_treasury.py}{111}{134}{The cash forecast: each account's P\&L, requirement, call or release each day, and the
survival horizon.}\label{lst:fm:treasury-and-funding:forecast}

\begin{definition}[Survival horizon]\label{def:fm:treasury-and-funding:horizon}
The \emph{survival horizon}\index{survival horizon} of a firm under a scenario is the number of days until its \omterm{def:fm:treasury-and-funding:cash}{unencumbered cash}, after meeting
every margin call and contractual outflow the scenario generates, first falls below zero, with no new funding.
\end{definition}

\section{Tutorial: three brokers and a bad week}\label{tut:fm:treasury-and-funding:tut}

\begin{tutorial}
\textbf{Goal.} Allocate the book across three brokers, then forecast thirty days of cash under a base case and a stressed week, with and without a
\omterm{def:fm:treasury-and-funding:lockup}{margin lock-up}. \textbf{End state:} \cref{tab:fm:treasury-and-funding:alloc} and the runway chart (\cref{fig:fm:treasury-and-funding:runway}).
\begin{tutsteps}
\item \textbf{The allocation.} \texttt{fm\_treasury.allocations()} builds the five allocations; \texttt{firm.treasury.optimise} solves the linear
programme with scipy's HiGHS solver.
\item \textbf{The accounts.} Each prime-broker account starts with its margin plus a 10\% cushion; the firm has \$40 million of \omterm{def:fm:treasury-and-funding:cash}{unencumbered cash}.
\item \textbf{The stress.} For five days the longs fall 2\% a day, the shorts 1\% and the index 2\%; then five flat days and a slow recovery. The
brokers raise their \omterm{def:fm:treasury-and-funding:house}{house margin} by a fifth a day, to double by day 5, hold it to day 15, then ease to one and a half times.
\item \textbf{The runway.} \texttt{fm\_treasury.run(stress, notice)} forecasts the cash with a lock-up of \texttt{notice} days (0 for none).
\end{tutsteps}
\end{tutorial}

\begin{omfigure}
\begin{tikzpicture}
\begin{axis}[width=10.5cm,height=5.4cm,xlabel={\small day},ylabel={\small unencumbered cash (\$ million)},xmin=0,xmax=30,ymin=-35,ymax=50,
  tick label style={font=\footnotesize},grid=major,grid style={gray!20},table/col sep=comma,
  legend cell align=left,legend style={font=\scriptsize,at={(0.5,-0.3)},anchor=north,/tikz/every even column/.append style={column sep=8pt}},legend columns=2]
\addplot[gray,thick] table[x=day,y=base]{figdata/desk/14-treasury-and-funding/runway.csv};
\addplot[omThm,thick,mark=*,mark size=1pt] table[x=day,y=lockup]{figdata/desk/14-treasury-and-funding/runway.csv};
\addplot[omDef,very thick] table[x=day,y=stress]{figdata/desk/14-treasury-and-funding/runway.csv};
\addplot[omProp,thick,dashed] table[x=day,y=one]{figdata/desk/14-treasury-and-funding/runway.csv};
\draw[black,thin] (axis cs:0,0) -- (axis cs:30,0);
\legend{base case,{stress, 30-day lock-up},{stress, no lock-up},{stress, all at broker B}}
\end{axis}
\end{tikzpicture}

\omcaption{Thirty days of \omterm{def:fm:treasury-and-funding:cash}{unencumbered cash}. Without a lock-up the optimised book runs out of cash on day 5, and the same book at a single broker on
day 3; with a thirty-day lock-up it keeps \$26.7 million at the worst. Data: \texttt{fm\_treasury.run}.}\label{fig:fm:treasury-and-funding:runway}
\end{omfigure}

The stressed week costs the book \$15.5 million of P\&L; the brokers' margin increases cost more. Without a lock-up the total requirement rises from
\$43.0 million on day 1 to \$71.8 million on day 5 (\cref{fig:fm:treasury-and-funding:req}), the firm pays \$57.9 million of calls over the month,
and its \$40 million runs out on day 5: it needed a buffer \$8.1 million larger. The same book at broker B alone runs out on day 3 and needed \$29.8
million more. With a thirty-day lock-up the requirement stays near \$37 million, the calls total \$23.0 million, and the firm ends the worst day with
\$26.7 million. A lock-up shorter than the stress only postpones the day: five days of notice move it to day 6, ten days to day 11. The clearing
broker's margin on the future hardly moves, from \$2.02 million to \$2.13 million, because its model responds to volatility with a lag; the large
term is the brokers' discretion.

\begin{omfigure}
\begin{tikzpicture}
\begin{axis}[width=10.5cm,height=4.8cm,xlabel={\small day},ylabel={\small total requirement (\$ million)},xmin=1,xmax=30,ymin=0,ymax=80,
  tick label style={font=\footnotesize},grid=major,grid style={gray!20},table/col sep=comma,
  legend cell align=left,legend style={font=\scriptsize,at={(0.5,-0.32)},anchor=north,/tikz/every even column/.append style={column sep=8pt}},legend columns=2]
\addplot[omDef,very thick,const plot] table[x=day,y=nolock]{figdata/desk/14-treasury-and-funding/req.csv};
\addplot[omThm,thick,dashed,const plot] table[x=day,y=lock]{figdata/desk/14-treasury-and-funding/req.csv};
\legend{no lock-up,30-day lock-up}
\end{axis}
\end{tikzpicture}

\omcaption{The total margin requirement across the four accounts in the stressed month. The gap between the lines is the brokers' increase in
\omterm{def:fm:treasury-and-funding:house}{house margin}, which a lock-up defers. Data: \texttt{fm\_treasury.paths}.}\label{fig:fm:treasury-and-funding:req}
\end{omfigure}

\section{Funding in stress}

Stress turns every lever of the forecast at once: positions lose money, requirements rise, brokers stop releasing excess, and funding sources
that were cheap become unavailable. The margin spiral of Book 3 is the market-wide version: forced sales to meet calls move prices, which raise
calls elsewhere. The firm's defences are the buffer, the terms and the diversity of its funding.

Diversity has three parts. Several brokers, so that one broker's decision moves only part of the requirement; several sources of cash, so that a
money market fund that gates or a bank that declines to lend is not the only route; and assets whose value does not fall with the book's, since the
collateral the firm posts in stress is worth what it will fetch that day. A committed credit line is the most expensive of the three in calm
markets, since the firm pays a fee on money it does not use, and the one whose value is highest on the day a call arrives that nothing else can
meet (\cref{fig:fm:treasury-and-funding:map}).

\begin{dated}{2026-09}{International guidance on margin preparedness}\label{dat:fm:treasury-and-funding:fsb}
After the March 2020 episode, the BCBS, CPMI and IOSCO reviewed margining practices (final report, September 2022), and the Financial Stability
Board published eight recommendations for non-bank market participants on liquidity preparedness for margin and collateral calls (final report, 10
December 2024): include the liquidity risk of margin and collateral calls in their risk management and governance, with contingency funding plans;
stress test their liquidity for such calls under extreme but plausible scenarios; and hold sufficient cash and readily available, diverse liquid
assets, with appropriate collateral arrangements and regular contact with their counterparties.
\end{dated}

\begin{method}[Setting the liquidity buffer]\label{met:fm:treasury-and-funding:buffer}
\begin{enumerate}
\item Map every account: what it holds, its requirement rule, whether excess is released, and the notice before its terms can change.
\item Build stress scenarios that move P\&L and requirements together, including brokers' discretionary increases and a counterparty that stops
releasing excess.
\item Forecast cash daily over the horizon of the longest notice period; the buffer is the largest cumulative shortfall.
\item Negotiate lock-ups for the terms that move most in the scenarios, and spread financing so that no single broker's decision exhausts the buffer.
\item Report the \omterm{def:fm:treasury-and-funding:horizon}{survival horizon} under each scenario to the \omterm{def:fm:the-risk-management-function:cro}{risk committee} monthly and after every material change in terms or positions.
\end{enumerate}
\end{method}

Book 1 tells how a family office's margin terms at several banks, and the banks' different reactions when its positions fell, turned into
large losses for some of them. For the borrower the lesson is this chapter's: the terms that matter are the ones that can change
in a week.

\section{Build: the treasury engine}\label{bld:fm:treasury-and-funding:treasury}

\begin{build}
\bfield{Purpose} \omterm{def:fm:treasury-and-funding:house}{House margin} as data, the allocation of a book across brokers, and the daily cash and collateral forecast with its \omterm{def:fm:treasury-and-funding:horizon}{survival
horizon}.
\bfield{Interface} \texttt{firm.treasury}: \texttt{Broker}, \texttt{margin}, \texttt{financing}, \texttt{annual\_cost}, \texttt{optimise};
\texttt{fcm\_margin} (on \texttt{firm.initmargin}), \texttt{house\_path}, \texttt{forecast}, \texttt{survival\_horizon}, \texttt{buffer\_needed}.
\bfield{Rules} Allocations are fractions that sum to one per position; the optimiser's cost never exceeds any single broker's; a call is paid the
day it arises; a lock-up holds the house multiplier at one for its notice period.
\bfield{Acceptance tests} \texttt{code/firm/treasury/tests/}: margin and financing on hand numbers; the optimiser against every single broker; a
three-day forecast with a retained excess; the clearing margin's scaling.
\bfield{Stretch} Collateral eligibility and haircuts; a revolving credit line with a commitment fee; the stress generated by firm.margin's scenario
engine for the clearing broker.
\end{build}

\begin{omsources}
\item BCBS, CPMI and IOSCO, \emph{Review of margining practices}, September 2022, and press release of 29 September 2022.
\item Financial Stability Board, \emph{Liquidity Preparedness for Margin and Collateral Calls: Final report}, 10 December 2024.
\item 12 CFR 220.12; FINRA Rule 4210.
\end{omsources}

\section{Exercises}

\begin{exercise}[$\star$]\label{exo:fm:treasury-and-funding:1}
Compute the \omterm{def:fm:treasury-and-funding:house}{house margin} of a single \$120 million long at each of the chapter's three brokers.
\end{exercise}

\begin{exercise}[$\star$]\label{exo:fm:treasury-and-funding:2}
What is the Regulation T initial margin on a \$10 million long in a margin equity security, and on a \$10 million short sale?
\end{exercise}

\begin{exercise}[$\star$]\label{exo:fm:treasury-and-funding:3}
The book's margin falls from \$72.0 million to \$33.9 million. At a 5\% cost of cash, what is that worth a year before financing?
\end{exercise}

\begin{exercise}[$\star\star$]\label{exo:fm:treasury-and-funding:4}
Show that holding the whole book at one broker is a feasible point of the linear programme, and explain why the optimised cost cannot exceed it.
\end{exercise}

\begin{exercise}[$\star\star$]\label{exo:fm:treasury-and-funding:5}
Why does broker B's rule make it the cheapest single broker for this book, although its margin is not the lowest?
\end{exercise}

\begin{exercise}[$\star\star$]\label{exo:fm:treasury-and-funding:6}
A lock-up of ten days is offered against one of thirty. On the chapter's stress, what does each buy?
\end{exercise}

\begin{exercise}[$\star\star\star$]\label{exo:fm:treasury-and-funding:7}
\emph{Coding.} Rerun the stress with starting cash of \$50 and \$60 million, for the optimised book and for the book at broker B. When does each
run out?
\end{exercise}

\begin{exercise}[$\star\star\star$]\label{exo:fm:treasury-and-funding:8}
\emph{Find the flaw.} ``We moved everything to the broker with the lowest margin; that frees cash, so our liquidity is better.''
\end{exercise}

\section{Problem: Three Brokers and a Bad Week}

\begin{problem}[{Weekend problem --- three brokers and a bad week}]\label{pb:fm:treasury-and-funding:1}
A new treasurer must decide where the firm holds its book and how much cash it keeps free.

\medskip\noindent\textbf{Part I --- The function.}
\begin{enumerate}
\item Define \omterm{def:fm:treasury-and-funding:cash}{unencumbered cash} and a \omterm{def:fm:treasury-and-funding:cash}{liquidity buffer}.
\item What happened to margin in March 2020?
\item Define \omterm{def:fm:treasury-and-funding:house}{house margin} and a \omterm{def:fm:treasury-and-funding:lockup}{margin lock-up}.
\item State the US regulatory floors for customer margin.
\end{enumerate}

\medskip\noindent\textbf{Part II --- The allocation.}
\begin{enumerate}[resume]
\item Define \omterm{def:fm:treasury-and-funding:optim}{margin optimisation}.
\item State and prove \cref{prop:fm:treasury-and-funding:lp}.
\item Give the margin and annual cost of the five allocations.
\item How does the optimiser use each broker's rule?
\item What does the optimiser not see?
\end{enumerate}

\medskip\noindent\textbf{Part III --- The forecast.}
\begin{enumerate}[resume]
\item Define the \omterm{def:fm:treasury-and-funding:horizon}{survival horizon}.
\item Describe the stress scenario and the P\&L it causes.
\item Give the total requirement on days 1 and 5, with and without lock-up, and the clearing broker's margin.
\item Give the \omterm{def:fm:treasury-and-funding:horizon}{survival horizons} and the buffers needed, for the optimised book and for the book at one broker.
\item What does a lock-up shorter than the stress buy?
\end{enumerate}

\medskip\noindent\textbf{Part IV --- The decision.}
\begin{enumerate}[resume]
\item What do the international recommendations of 2024 ask of non-bank firms?
\item How would you set the buffer?
\item Which is worth more on this stress: the optimisation or the lock-up?
\item What would you change in the brokers' agreements?
\item State the \emph{named result}: the margin saved by optimising across three brokers against using one, and the \omterm{def:fm:treasury-and-funding:horizon}{survival horizon} in the stress
week with and without \omterm{def:fm:treasury-and-funding:lockup}{margin lock-up}.
\item In two sentences, write the treasury policy.
\end{enumerate}
\end{problem}

\section{Interview questions}

\begin{interviewq}[$\star$ \iqroles{trader}]\label{iq:fm:treasury-and-funding:1}
What is \omterm{def:fm:treasury-and-funding:house}{house margin}, and why is it usually above the regulatory minimum?
\end{interviewq}

\begin{interviewq}[$\star$ \iqroles{risk}]\label{iq:fm:treasury-and-funding:2}
What is a \omterm{def:fm:treasury-and-funding:horizon}{survival horizon}, and how would you compute it?
\end{interviewq}

\begin{interviewq}[$\star\star$ \iqroles{researcher, developer}]\label{iq:fm:treasury-and-funding:3}
Formulate the allocation of positions across brokers with gross, net and concentration margin terms as an optimisation problem.
\end{interviewq}

\begin{interviewq}[$\star\star$ \iqroles{risk}]\label{iq:fm:treasury-and-funding:4}
Your prime broker doubles its margin overnight. What do you check, and what do you do?
\end{interviewq}

\begin{interviewq}[$\star\star$ \iqroles{trader, risk}]\label{iq:fm:treasury-and-funding:5}
Why might a firm pay more margin than it needs to?
\end{interviewq}

\begin{interviewq}[$\star\star\star$ \iqroles{risk, researcher}]\label{iq:fm:treasury-and-funding:6}
Design a liquidity stress test for a hedge fund with three prime brokers and a futures account.
\end{interviewq}
